\clearpage
\hypertarget{gaussix-sources}{}
\section*{Quellen und Textanmerkungen}

\subsection*{Historische Ausgabe und Mitwirkende}

Der Übersetzung liegt die originalsprachliche KI-Transkription von Band IX vom 11. November 2026 zugrunde, \texttt{gauss-werke-ix-1-528-latex.zip}.

\noindent Carl Friedrich Gauss, \textit{Werke. Neunter Band}. Herausgegeben von der Königlichen Gesellschaft der Wissenschaften zu Göttingen. In Kommission bei B. G. Teubner, Leipzig, 1903. Felix Klein war Mitherausgeber der Werkausgabe und zeichnete die abschließenden \textit{Bemerkungen zum neunten Bande} als F. Klein. \href{https://gdz.sub.uni-goettingen.de/id/PPN23601515X}{Katalog des Göttinger Digitalisierungszentrums}. \href{https://lbezone.hkust.edu.hk/rse/?p=13258}{Katalog der HKUST Library}.

Carl Friedrich Gauss’ mathematische Schriften, Manuskripte und Briefe bilden den Hauptteil des Bandes. Klein nennt auf S.~521 L. Krüger als Bearbeiter des neunten Bandes und A. Börsch als seinen Mitarbeiter bei der Sichtung des Materials. M. Brendel leitete die allgemeine Redaktion wie schon bei Band VIII. Die historischen Bemerkungen tragen ihre ursprünglichen Unterschriften, darunter die von Krüger und die gemeinsamen Unterschriften von Krüger und Börsch. Bei den Stationsbeobachtungen ist Hartmann genannt. Die Einleitung zum ersten veröffentlichten Bericht über die hannoversche Triangulation stammt von H. C. Schumacher. Die Bemerkungen zu den Stationsausgleichungen zitieren aus Goldschmidts \textit{Geodätischer Calcül nach Gauss}. Leutnant J. Gauss entwarf die hannoversche Dreieckskarte. Krügers begleitende historische Anmerkung bezeichnet die redaktionellen Berichtigungen an der Karte.

Die historischen Herausgeber kennzeichnen ihre Einschübe mit eckigen Klammern oder in den Bemerkungen und setzen Auszüge aus Notizen und Briefen anderer Verfasser in geschweifte Klammern. Klein erläutert diese Konventionen auf S.~521.

Schrift und Anordnung des Titelblatts folgen dem Göttinger Exemplar von 1903. Die unsignierte Vignette ist bereits auf dem \href{https://gdz.sub.uni-goettingen.de/id/PPN235993352}{Titelblatt von \textit{Werke} I (1863)} nachgewiesen. Die hier abgebildete Vignette stammt aus dem zweiten Abdruck von \textit{Werke} I (1870), \href{https://commons.wikimedia.org/wiki/File:Gauss,_Carl_Friedrich_-_Vignette_-_von_Werke_(1870)_(page_9_crop).jpg}{Wikimedia Commons, Version 1179960952}. \href{https://creativecommons.org/publicdomain/mark/1.0/}{Public Domain Mark 1.0}. Der historische Text, die Formeln, Tabellen und mathematischen Zeichnungen liegen in editierbarem LaTeX mit modernen Schriften und Abständen vor. Die Zeichnungen folgen den gedruckten Beschriftungen und Verbindungen mit den unten angegebenen Annäherungen der Geometrie. Der historische Vorspann umfasst seine drei leeren Rückseiten. Der moderne Digitalisierungsumschlag, Bibliotheksstempel und handschriftliche Zugangsvermerke gehören zu einzelnen Exemplaren.

\subsection*{Inhalt des neunten Bandes}

\noindent Die Seitenzahlen beziehen sich auf die Ausgabe von 1903 und schließen die nach ihrer Stellung mitgezählten unnummerierten Seiten ein. Auf S.~528 folgen das letzte unnummerierte Leerblatt und die Verlagsanzeige. Das historische Inhaltsverzeichnis behält seine ursprünglichen Seitenverweise.

\paragraph{Bestimmung des Breitenunterschiedes zwischen Göttingen und Altona (S.~1--64).}
Die \textit{Bestimmung des Breitenunterschiedes zwischen den Sternwarten von Göttingen und Altona durch Beobachtungen am Ramsdenschen Zenithsector} enthält Titel, Inhaltsverzeichnis und Einleitung, das Verzeichnis von 43 Sternen und die Beobachtungsreihen. Die Resultate umfassen Gewichtsformeln, Kollimationsfehler, Residuen, statistische Reduktionen, direkte und reflektierte Beobachtungen und Sternvergleiche. Die Beobachtungen auf dem Seeberg umfassen sechzig Ergebnisse. Der \textit{Zusatz} nennt fünfundzwanzig Vermessungsorte, die \textit{Anzeige} enthält historische Bemerkungen.

\paragraph{Erdellipsoid und geodätische Linie (S.~65--98).}
Auf das Titelblatt mit seiner leeren Rückseite folgen drei Nachlassartikel über das Ellipsoid mit ihren numerischen Rechnungen und Bemerkungen sowie die Artikel [1]--[13] zur geodätischen Linie. Sie enthalten Hilfstafeln, zwei Tabellen mit je 61 Zeilen, ausgeführte Rechnungen, mathematische Zeichnungen und die historischen Bemerkungen auf S.~96--98.

\paragraph{Polhöhe und sphärische Trigonometrie (S.~99--104).}
Der Brief an J. J. Baeyer über die Änderung der Polhöhe mit der Höhe steht auf S.~99--102 mit Zeichnung und historischen Bemerkungen. Zwei weitere Nachlassnotizen behandeln die Winkel sphärischer Dreiecke und Punkte auf einem Großkreis. Dazu kommen eine Zeichnung und die von Krüger und Börsch gezeichneten Bemerkungen.

\paragraph{Konforme Doppelprojektion (S.~105--122).}
Die \textit{Conforme Doppelprojection des Sphäroids auf die Kugel und die Ebene} beginnt mit Titelblatt und leerer Rückseite. Die Artikel [1]--[10] zur Abbildung des Ellipsoids auf die Kugel enthalten eine Zahlentafel und historische Bemerkungen. Darauf folgen die stereografische Projektion der Kugel auf die Ebene und die zugehörigen Bemerkungen.

\paragraph{Mercator-Projektion und stereografische Darstellung des Sphäroids (S.~123--134).}
Zu den Artikeln [1]--[8] über die Mercator-Projektion gehören eine Zahlentafel und die von Krüger und Börsch gezeichneten Bemerkungen. Die stereografische Darstellung des Sphäroids in der Ebene umfasst die Formeln und historischen Bemerkungen auf S.~133--134.

\paragraph{Konforme Abbildung des Sphäroids auf den Kegelmantel (S.~135--140).}
Auf Titelblatt und leere Rückseite folgen der Nachlasstext mit Formeln und zwei Zahlentafeln und die historischen Bemerkungen.

\paragraph{Konforme Abbildung des Sphäroids in der Ebene (S.~141--204).}
Auf Titelblatt und leere Rückseite folgen die Artikel [1]--[32] aus dem Nachlass und die historischen Bemerkungen. Dazu gehören Koordinatentransformationen, logarithmische und Krümmungsformeln, Zahlentafeln und die ausgeführte Rechnung für Varel.

\paragraph{Über die Formeln für die hannoversche Landesvermessung (S.~205--218).}
Die Briefe an Schumacher sind auf den 18. April 1830, 30. April 1830, 17. Mai 1831, 25. Juni 1831 und 9. Dezember 1838 datiert. Sie enthalten Hilfstafeln und Rechnungen für Neuwerk. Die historischen Bemerkungen sind von Krüger und Börsch gezeichnet.

\paragraph{Trigonometrische Punktbestimmung (S.~219--242).}
Auf Titelblatt und leere Rückseite folgen die Artikel [1]--[9] und die historischen Bemerkungen mit Zeichnungen, Tabellen, dem Zitat aus Goldschmidt und der Nennung Krügers.

\paragraph{Ausgleichung einfacher Figuren (S.~243--260).}
Auf Titelblatt und leere Rückseite folgen die Artikel [1]--[7], die Briefe vom 11. Februar 1824 und 19. Januar 1840 und die von Krüger gezeichneten Bemerkungen. Die Artikel enthalten Zeichnungen, Formeln und Tabellen.

\paragraph{Stationsausgleichungen (S.~261--293).}
Auf Titelblatt und leere Rückseite folgen die Beobachtungen und Normalgleichungen für Zeven, Brillit, Wilsede und Windberg mit der historischen Nennung Hartmanns. Darauf folgen die Beobachtungen am Breithorn vom 13.--16. September 1822, der Brief an Gerling vom 26. Dezember 1823 und die Briefe an Schumacher vom 22. Dezember 1827 und 7. Januar 1828. Auf die Tabellen, die Dreieckszeichnung und die historischen Anmerkungen folgen Krügers Bemerkungen, die aus Goldschmidts \textit{Geodätischer Calcül nach Gauss} zitieren.

\paragraph{Zur Netzausgleichung (S.~294--328).}
Auf Titelblatt und angrenzende Leerseiten folgen die Formeln für die Anzahl der Bedingungsgleichungen [1], die 33 Einträge des Verbindungsverzeichnisses und das Netz mit 75 Kanten [2] sowie die 51 Gruppen beobachteter Dreiecke mit ihren abschließenden Relationen [3]. Artikel [4] enthält 51 nummerierte Gleichungen. Artikel [5] umfasst die Einleitung und die zwölf ausgeführten Fälle I--XII. Artikel [6] enthält die Relationstabelle, die Korrekturausdrücke und zwölf Normalgleichungen. Artikel [7] enthält die 150 Richtungskorrekturen und den Schlusstext. Artikel [8] enthält 51 Dreiecksgruppen und die Verbindungstabelle Hamburg--Hohenhorn. Artikel [9] enthält 150 Zeilen mit Azimuten und Reduktionen. Artikel [10] umfasst Auszüge aus Gauss' Briefen an Olbers vom 2. November 1823, 6. Juli 1824 und 14. Mai 1826 sowie an Gerling vom 2. Mai 1837, 5. Juni 1838 und 14. November 1838. Dazu gehören die historischen Fußnoten, die Triangulationszeichnung mit sechs Punkten, die Zeichnung des Viereckskranzes und Krügers historische Bemerkungen.

\paragraph{Dreieckskranz um Oldenburg (S.~329--342).}
Auf Titelblatt, leere Rückseite und einleitenden Absatz folgen 84 vorbereitende Tabellenzeilen in 21 Stationsgruppen, ein Triangulationsnetz mit 21 nummerierten Knoten und 42 Kanten, die 22 Winkelsummenfehler, Beispiele zu den Bedingungen, 22 Normalgleichungen und Korrelaten und die Erörterung der Korrekturen. Die Dreiecksübersicht enthält 63 Zeilen mit Winkeln und Seitenlogarithmen. Die Koordinatentabelle umfasst 23 Koordinatenzeilen und zwei Differenzzeilen. Die zweite Ausgleichung enthält die Formeln mit sechs Richtungsrelationen, die 77 Koeffizientenpaare und deren Erläuterung. Krügers historische Bemerkungen zitieren Gauss' Brief an v.~Schrenck vom 3. August 1835. Die historische Ausgleichung selbst ist unvollendet.

\paragraph{Zur hannoverschen Triangulation (S.~343--434 und Faltkarte).}
Auf Titelblatt und leere Rückseite folgt der Briefwechsel. Abschnitt [1] umfasst acht Auszüge aus Briefen an Schumacher vom 5. Juli 1816 bis zum 14. Januar 1829 und drei historische Fußnoten. Abschnitt [2] umfasst acht Auszüge aus Briefen an Bessel vom 26. Dezember 1821, 15. November 1822, 5. November 1823, 20. November 1824, 12. März 1826, 20. November 1826, 1. April 1827 und 9. April 1830 mit drei historischen Fußnoten. Abschnitt [3] enthält den Brief an Bohnenberger vom 16. November 1823 mit seinen drei historischen Fußnoten. Abschnitt [4] enthält zwölf Auszüge aus Briefen an Olbers vom 13. Januar 1821, 18. April 1822, 4. Juli 1824, 8. Juli 1824, 19. Februar 1825, 25. Februar 1825, 9. Oktober 1825, 2. April 1826, 14. Januar 1827, 1. März 1827, 14. Juni 1830 und 2. September 1837. Dazu gehören drei historische Fußnoten, die Dreieckszeichnung Bremen--W.--Brüttendorf--Bottel und die 57-zeilige Tabelle \textit{Relative Höhen}. Abschnitt [5] enthält elf Auszüge aus Briefen an Gerling vom 5. Oktober 1821, 21. Februar 1822, 7. November 1822, 27. Juli 1823, 11. August 1823, 1. September 1823, 5. September 1823, 3. Oktober 1823, 19. Juli 1827, 12. September 1838 und 14. November 1838. Dazu gehören vier Zeichnungen zur Signalgebung mit Heliotropen, drei optische Zeichnungen, eine Koordinatentabelle, eine historische Fußnote und die von Krüger gezeichneten Bemerkungen.

Abschnitt [II.] enthält die veröffentlichten Berichte von 1821 und 1822, Schumachers Einleitung, die Koordinatentabelle und die historischen Bemerkungen. Abschnitt [III.1] enthält den Plan der Arbeit, seine Einleitung und den Anfang seines ersten Abschnitts. Abschnitt [III.2] enthält Auszüge aus den Berichten vom 7. Januar 1822, 31. Januar 1823, 16. Februar 1825, 21. November 1827 und 26. Juni 1828, den historischen Bericht vom 8. Februar 1838 und die Auszüge vom 5. Juli 1840 und Dezember 1844. Abschnitt [III.3] enthält die vollständige 33-zeilige Koordinatentabelle und Krügers historische Bemerkungen. Die editierbare Faltkarte zeigt die Dreiecksnetze, geografische Umrisse, die Legenden und den Maßstab.

\paragraph{Höhenmessungen (S.~435--458).}
Auf Titelblatt und leere Rückseite folgen die drei Texte zur Refraktion [1]--[3], darunter die Briefe an Schumacher vom 28. November 1824 und 27. Dezember 1846 und die historische redaktionelle Anmerkung. Der Nachlass zur terrestrischen Refraktion und Höhenausgleichung enthält 71 Fehlergleichungen, 31 Normalgleichungen, Stationskorrekturen und die Bemerkungen Krügers. Drei Tabellen begleiten die Anleitung zur barometrischen Höhenmessung.

\paragraph{Heliotrop (S.~459--484).}
Auf Titelblatt und leere Rückseite folgen der Artikel in den \textit{Göttingischen gelehrten Anzeigen} vom 9. August 1821 mit einer optischen Zeichnung und der am 26. Dezember 1821 eingereichte Beitrag für das \textit{Astronomische Jahrbuch}. Der Brief an Olbers ist auf den 1. Juli 1821 datiert. Die vier Briefe an Schumacher stammen vom 11. Juli 1821, 8. November 1821, 10. Juli 1822 und 15. Januar 1827, mit Anmerkungen und Zeichnung. Der Artikel über die Justierung des Heliotrops enthält eine Tafel mit dreizehn Figuren. Zwei Nachlassnotizen enthalten Formeln und eine geometrische Zeichnung. Krügers historische Bemerkungen enthalten eine Zeichnung der ersten Konstruktion.

\paragraph{Messungsfehler (S.~485--506).}
Auf Titelblatt und leere Rückseite folgen der Text über Vermessungsinstrumente, die Briefe über Messungsfehler an Olbers (Juli 1825), Schumacher (14. August 1825), Bessel (29. Oktober 1843 und 15. August 1844) und Gerling (17. April 1844) und die Rechnungen aus den Beobachtungen. Die historischen Bemerkungen sind von Krüger gezeichnet.

\paragraph{Reduktion schiefer Winkel auf den Horizont. Repetitionsbeobachtungen. (S.~507--520).}
Auf Titelblatt und leere Rückseite folgen die drei Notizen über die Reduktion schiefer Winkel auf den Horizont, ein ausgeführtes Beispiel, das Rechnungsschema und die historischen Bemerkungen. Zu den Repetitionsbeobachtungen gehören der Brief an Bessel vom 27. Januar 1819, die Rekursions- und Fehlerformeln sowie die von Krüger gezeichneten historischen Bemerkungen.

\paragraph{Bemerkungen zum neunten Band und Inhalt (S.~521--528).}
F. Klein zeichnete die historischen Bemerkungen zum neunten Band. Das historische Inhaltsverzeichnis, das letzte Leerblatt und die Verlagsanzeige schließen den Band ab.

\subsection*{Druckexemplare und Karte}

Textgrundlage ist der vom Göttinger Digitalisierungszentrum digitalisierte
historische Druck von Carl Friedrich Gauss’ \emph{Werken}, Band IX
(PPN23601515X). \href{https://gdz.sub.uni-goettingen.de/download/pdf/PPN23601515X/PPN23601515X.pdf}{Digitalisat des Göttinger Digitalisierungszentrums}.

Der historische Umfang umfasst die gedruckten Seiten 1--528 mit der
unpaginierten Faltkarte zur hannoverschen Triangulation. Dazu gehören die
Titelblätter zu \emph{Conforme Doppelprojection des Sphäroids auf die Kugel
und die Ebene}, \emph{Conforme Übertragung des Sphäroids auf den
Kegelmantel}, \emph{Conforme Abbildung des Sphäroids in der Ebene},
\emph{Trigonometrische Punktbestimmung}, \emph{Ausgleichung einfacher
Figuren}, \emph{Stationsausgleichungen}, \emph{Zur Netzausgleichung},
\emph{Dreieckskranz um Oldenburg}, \emph{Zur Hannoverschen Triangulation}, \emph{Höhenmessungen}, \emph{Heliotrop},
\emph{Messungsfehler} und \emph{Reduction schiefer Winkel}. Zum Band
gehören ferner der historische Vorspann mit drei leeren Rückseiten, die
Leerseiten an den Positionen 2, 4, 66, 106, 136, 142, 220, 244, 262, 294,
296, 330, 344, 436, 460, 486, 506 und 508 sowie das unpaginierte leere
Schlussblatt und die Verlagsanzeige.

Das \href{https://archive.org/details/carlfriedrichga01gausgoog9}{Digitalisat von Google Books im Internet Archive} dient als Vergleichsexemplar.

Das Exemplar der HKUST Library macht schwach gedruckte Klammern,
kleine Schriftzeichen und weitere Einzelheiten auf S.~53--54, 57--58,
60, 64, 69--72, 76--77, 80, 82, 86, 98, 102, 109--112, 119, 125, 129,
132, 138, 144, 152--162, 171, 195--199, 201--202 und 212 deutlicher
lesbar. Es stützt die Lesung der Verbindungsliste und des Netzes auf
S.~298--299 sowie der beiden Tabellenüberschriften
\(\tfrac13\,\text{Excess}\) auf S.~304--305. Auf S.~441 klärt es den
beschädigten Buchstaben in \emph{20-maliger}.
\href{https://lbezone.hkust.edu.hk/rse/?p=13258}{Vergleichsexemplar der HKUST Library}.
\href{https://lbezone.hkust.edu.hk/obj/4/o/b761099/v9/pg00220.jpg}{S.~212},
\href{https://lbezone.hkust.edu.hk/obj/4/o/b761099/v9/pg00312.jpg}{S.~304},
\href{https://lbezone.hkust.edu.hk/obj/4/o/b761099/v9/pg00313.jpg}{S.~305}
und
\href{https://lbezone.hkust.edu.hk/obj/4/o/b761099/v9/pg00451.jpg}{S.~441
im Exemplar der HKUST Library}.

Die auseinandergefaltete Karte \textit{Übersicht der gemessenen Dreieckssysteme} folgt
dem gesonderten Digitalisat der SUB Göttingen. Im Göttinger
Banddigitalisat fehlt die Karte. Das Google-Digitalisat zeigt nur
gefaltete Streifen. Geometrie, Strichstärken und der Versatz der
Farbbänder geben die historische Zeichnung näherungsweise wieder.
Bei Pugelatz und Fuhrenberg sowie am Strahl zum Inselsberg lassen sich
die Mittellinien überdruckter Konturen stellenweise nicht voneinander
unterscheiden. Die verdeckten Verläufe bleiben offen.
\href{https://webdoc.sub.gwdg.de/ebook/e/2005/gausscd/html/kapitel_landv_karte.htm}{Karte der hannoverschen Triangulation, SUB Göttingen}.

Der zitierte Brief von Gauss an v.~Schrenck vom 3.~August 1835
erschien auch im \emph{General-Bericht über die mitteleuropäische
Gradmessung für das Jahr 1865}, Berlin, Georg Reimer, 1866, S.~23--24.
\href{https://books.google.com/books?id=VBNDAAAAcAAJ\&pg=PA24}{Früherer Druck des Briefes, Google Books}.

\subsection*{Lesarten der Druckexemplare}

Die zitierten Schreibweisen stammen aus den historischen Exemplaren. Die Anmerkungen erläutern auch die Wiedergabe der mathematischen Zeichen und Zeichnungen in dieser Ausgabe.

Die beiden tief stehenden Zeichen im Kolumnentitel auf S.~6 und das
Zeichen nach 54,24 in der Beobachtungstabelle auf S.~18 folgen dem
Göttinger Exemplar. Ihre Funktion ist unklar. Die Sternbezeichnung
auf S.~9 verwendet dieselbe Kappa-Variante wie auf S.~51. Der
Doppelpunkt bei \emph{Altona} auf S.~12 und \emph{Jun} ohne Punkt
auf S.~15 folgen dem Göttinger Exemplar.

In der Gewichtsformel auf S.~30 steht nach \(\beta'\) ein einzelner
Punkt, dessen Funktion unklar ist. Die abweichenden Breitenwerte auf
S.~38 und 39 stehen so im Druck. Die Bogensignatur auf S.~35 ist im
Google-Digitalisat und im Exemplar der HKUST Library lesbar.

Der einzelne Punkt vor \emph{den} auf S.~51 folgt dem Göttinger
Exemplar. Seine Funktion ist unklar. Die Bogensignatur auf dieser
Seite folgt dem Google-Digitalisat. Das gedrungene kreuzförmige
griechische Zeichen in Sternzeile 8 wird mit der Kappa-Variante
wiedergegeben. Die genaue historische Zeichenform bleibt unsicher.
Die Zahlen und Ortsnamenschreibungen in den Seeberger Tabellen und
den Vermessungstabellen folgen dem Druck, einschließlich der
unreduzierten Sekundenangaben.

Schwach gedruckte Begrenzungszeichen von Zitaten auf S.~64 und 71,
kleine Ziffern auf S.~69, Klammern und die äußere Quadrierung auf
S.~70 sowie eine schwach gedruckte Klammer auf S.~72 folgen dem
deutlicheren Exemplar der HKUST Library. Der einzelne Punkt vor dem
letzten Logarithmus auf S.~69 folgt dem Göttinger Exemplar. Seine
Funktion ist unklar. Die Hilfstabelle auf S.~77 enthält die leeren
Zellen und abgekürzten Werte des Drucks. Die Zeichnungen auf S.~73
und 75 geben die beschrifteten Verbindungen mit näherungsweiser
Geometrie wieder.

Die im Göttinger Digitalisat abgeschnittene rechte Tabellenspalte
auf S.~82 folgt dem Exemplar der HKUST Library. Die Abplattungswerte
auf S.~82 und 84 weichen schon im Druck von den Werten auf S.~83
und 87 ab. Die einzelnen Punkte auf S.~84, 88, 93 und 94 und der
hoch stehende Punkt am Formelende auf S.~91 bleiben ohne zusätzliche
mathematische Deutung erhalten. Eine schwach gedruckte schließende
Klammer auf S.~86 und die vierte Potenz auf S.~98 folgen dem
deutlicheren Exemplar der HKUST Library. Die Zeichnungen auf S.~100
und 103 geben die Verbindungen, Beschriftungen und Kreuzungen des
Drucks mit näherungsweiser Geometrie wieder.

Die auf S.~102 zitierte Seitenspanne 30--32 und kleine
Schriftzeichen auf S.~109--112 folgen dem deutlicheren Exemplar der
HKUST Library. Der einzelne Punkt vor dem Nenner auf S.~113 folgt
dem Göttinger Exemplar. Seine Funktion ist unklar. Der senkrechte
Strich nach \emph{cos} auf S.~119 ist im Göttinger, im Google- und
im HKUST-Exemplar vorhanden. Seine Funktion bleibt unklar. Nicht
eindeutig als Schriftzeichen erkennbare Spuren am Rand auf S.~112,
117--118 und 120--121 sind im Satz ausgespart.

Der tief stehende Punkt unter \mbox{2345.09} auf S.~126 ist getrennt
von den Ziffern wiedergegeben. Seine Funktion ist unklar. Der kurze
Strich bei \mbox{8504.24} lässt sich nicht eindeutig einordnen.
Eine schwach gedruckte schließende Klammer auf S.~129 folgt dem
deutlicheren Exemplar der HKUST Library. Dieses klärt auch den Nenner
45360 und die Klammern auf S.~132 sowie die Reihenentwicklung auf
S.~138. Die Zwischenräume in \mbox{6,709 8436 473} und
\mbox{1180 553} auf S.~139 folgen dem Göttinger Exemplar. Das
Google-Digitalisat zeigt an diesen Stellen Punkte. Der einzelne
Punkt vor 860 in der unteren Tabelle folgt dem Göttinger Exemplar
und fehlt im Google-Digitalisat. Die Funktion der einzelnen Punkte
auf S.~138--140 ist unklar. Nicht eindeutig als Schriftzeichen
erkennbare Spuren am Rand auf S.~132, 137 und 140 bleiben außerhalb
der mathematischen Notation.

Die schwach gedruckten schließenden Klammern auf S.~144 folgen dem
deutlicheren Exemplar der HKUST Library. Die Klammern in der ersten
Formel mit \emph{log hyp D} auf S.~147 folgen dem Google-Digitalisat.
Der Punkt nach 89,437 auf S.~160 folgt dem Göttinger Exemplar und
fehlt in den Vergleichsexemplaren. Die beiden Zeichen vor 109 auf
S.~161 und das hoch stehende, an ein Anführungszeichen erinnernde
Zeichen vor 89 auf S.~162 haben eine unklare Funktion. Der Tabellenwert
788 auf S.~162 und der Nenner 5040 auf S.~171 folgen den deutlicheren
Vergleichsexemplaren. Die einzelnen Punkte auf S.~163, 167 und
169--171 folgen dem Göttinger Exemplar. Ihre Funktion ist unklar.
Nicht eindeutig als Schriftzeichen erkennbare Spuren am Rand auf
S.~144, 147--150 und 166 bleiben außerhalb der mathematischen
Notation.

Auf S.~195 steht nach \emph{zugefügt worden} im Göttinger Exemplar
ein Komma, in den Google- und HKUST-Exemplaren ein Semikolon. Nach
der ersten abgesetzten Angabe \mbox{3,418 9411.803} auf S.~197 steht
im Göttinger Exemplar ein Punkt und in den Vergleichsexemplaren ein
Komma. Der klein gedruckte Koeffizient auf S.~195 und die eng gesetzten
Zeichen auf S.~196--199 und 201--202 sind im Exemplar der HKUST
Library deutlicher zu erkennen. Die vereinzelten Spuren auf S.~192
und 198 sowie im ansonsten leeren Gradfeld der Tabelle auf
S.~210 lassen sich nicht als Schriftzeichen einordnen. Die einzelnen
Punkte im Rechenfeld auf S.~203, nach \emph{Es} und neben der ersten
Längengleichung auf S.~208 sowie über dem Minuszeichen in Formel 6)
auf S.~209 folgen dem Göttinger Exemplar. Ihre Funktion ist unklar.

Auf S.~212 folgt der Exponent 2 im Nenner von \(F\) dem deutlicheren
Exemplar der HKUST Library. Die beiden Zeichen zwischen 434 und
2945 in \(k\), eines oberhalb und eines unterhalb, folgen dem
Göttinger Exemplar. Ihre Funktion bleibt unklar. Der einzelne kurze
Strich nach dem auf \emph{gebrauchen} folgenden Komma auf S.~225
folgt dem Göttinger Exemplar und fehlt im Google-Digitalisat. Ob
der Strich zum Text gehört, ist unklar.

Auf S.~233 folgt auf \emph{scheint} im Göttinger Exemplar ein
Semikolon, im Vergleichsexemplar ein Komma. Die Funktion des hoch
stehenden kommaähnlichen Zeichens beim Titel auf S.~239 ist
unklar. Auf S.~242 folgt auf \emph{sein} im Göttinger Exemplar ein
Punkt, im Vergleichsexemplar ein Doppelpunkt. Der Punkt vor
\emph{Rechnungen} und der hoch stehende Punkt vor \emph{eine} auf
S.~249 bleiben ohne Deutung erhalten. Auf S.~250 folgt der im
Google-Digitalisat verdeckte Text dem Göttinger Exemplar. Auf
S.~270 ist der rechte Teil des Systems im Göttinger Digitalisat
abgeschnitten. Das vollständige Google-Digitalisat zeigt diesen
Teil und seine Anordnung im Querformat. Die im Google-Digitalisat
unleserlichen Stellen auf S.~254 folgen dem Göttinger Exemplar.
Die einzelnen Zeichen auf S.~254--257, 260, 263--264 und 268--269
folgen dem Göttinger Exemplar. Ihre Funktion bleibt unklar. Die
wiederholten \(e\)-Überschriften auf S.~265 stehen so im Druck.

Die einzelnen Zeichen neben den Formeln auf S.~272--273, bei den
Seitenzahlen auf S.~274 und 276, beim Wert 253 auf S.~276, bei der
Fußnote auf S.~278 sowie an zwei Stellen im Fließtext und im
Formelsatz auf S.~281 folgen dem Göttinger Exemplar. Ihre Funktion
und ihr materieller Ursprung sind unklar. Der Schlussstrich auf
S.~277 ist im Google-Digitalisat sichtbar. Die Wiedergabe folgt hier
dem Göttinger Exemplar, in dem der Strich fehlt.

Im ersten Eintrag auf S.~286 steht nach \emph{I} im Göttinger
Exemplar ein Komma, im Vergleichsexemplar ein Punkt. Die tief
stehenden logarithmischen Zeichen auf S.~287 werden mit \(n\)
wiedergegeben. Ihre genaue gedruckte Form ist unsicher. Die einzelnen
Zeichen auf S.~292 und 297 sowie in der Wilsede-Zeile auf S.~298
folgen dem Göttinger Exemplar. Ihre Funktion ist unklar. Das erste
Füllzeichen in der Wilsede-Zeile ist im Göttinger Exemplar
kommaförmig, im HKUST-Exemplar punktförmig. Auf S.~300 steht in der
Summe für Dreieck 20 im Göttinger Exemplar \mbox{0.545}, im
Google-Digitalisat eine Zahl mit Komma. Der einzelne Punkt bei
Dreieck 8 und der Strich unter der mittleren Tabelle folgen ihrer
Position im Göttinger Exemplar. Auch der obere Punkt über dem
Komma bei Dreieck 48 auf S.~301 folgt diesem Exemplar. Einzelne
Zeichen auf S.~283, 293 und 301 sind nur in Vergleichsexemplaren
sichtbar. Die Wiedergabe folgt an diesen Stellen dem Göttinger
Exemplar. Spuren am Rand auf S.~282--285 lassen sich nicht als
zusätzliche Schriftzeichen einordnen. Das Netz auf S.~299 gibt die
Verbindungen und Beschriftungen des Drucks mit näherungsweiser
Geometrie wieder.

Auf S.~444 liest das Göttinger Exemplar \emph{Nach diesem}, bezogen auf das Handbuch. Der folgende Verweis auf dieses Handbuch steht in runden Klammern.

Die Bogensignatur auf S.~465 und der Schlussstrich auf S.~466 folgen
dem Google-Digitalisat. Beide fehlen im Göttinger Digitalisat.
Die einzelnen Punkte auf S.~473, 474 und 480 folgen dem Göttinger
Exemplar. Ihre Funktion ist unklar. Die Tafel auf S.~477 trägt die
gedruckten Bezeichnungen \emph{J} und \emph{c}, während der
Fließtext auf \emph{I} und \emph{CC} verweist.

Im Göttinger Exemplar steht auf S.~489 ein Semikolon nach \emph{Müller}, im Google-Digitalisat ein Komma. Der hoch stehende Punkt auf S.~495 folgt dem Göttinger Exemplar und fehlt im Google-Digitalisat. Der abschließende Punkt nach \emph{Bremervörde} in der
Zeichnung auf S.~492 ist im Göttinger Exemplar undeutlich und im
Google-Digitalisat punktförmig. Die unregelmäßige Druckspur am
lesbaren Wort \emph{Ich} auf S.~495 lässt sich nicht als
zusätzliches Zeichen einordnen.

Die einzelnen Punkte vor \emph{Fernrohr} auf S.~501 und vor
\emph{Fehler} auf S.~515 folgen dem Göttinger Exemplar. Im
Google-Digitalisat fehlen sie. Ihre Funktion ist unklar. Die
Bogensignatur \emph{IX. / 64} auf S.~505 folgt dem Google-Digitalisat
und fehlt im Göttinger Digitalisat. Schwach gedruckte Klammern auf
S.~512 folgen dem besser lesbaren Vergleichsexemplar. Das
klammerartige Zeichen zwischen dem Gleichheitszeichen und der
Kubikwurzel in der Definition von \(\mu\) auf S.~513 wird als
schließende eckige Klammer \texttt{]} wiedergegeben. Seine Funktion ist
unklar. Unregelmäßige Spuren bei \emph{Aus} und beim letzten
Exponenten auf S.~516 lassen sich nicht als zusätzliche
Schriftzeichen einordnen.

Der Exponent 2 in der Eliminationsgleichung auf S.~519 und der
Exponent 4 im Zähler des reziproken Gewichts folgen dem Göttinger
Exemplar. Im Google-Digitalisat sind beide Hochziffern kaum zu
entziffern. Das Datum 18.~September 1822 auf S.~526 folgt
dem Göttinger Exemplar. Im Google-Digitalisat scheint der Tag als
19 gedruckt zu sein. Der einzelne Punkt nach \emph{Auszug} auf
S.~526 ist wiedergegeben. Spuren bei der Überschrift auf S.~521,
bei \emph{Azimuth} auf S.~524 und bei \emph{an Olbers} auf S.~526
lassen sich nicht als zusätzliche Schriftzeichen einordnen.

Auf S.~528 steht im Göttinger Exemplar \emph{Uber Messungsfehler} ohne Umlaut in \emph{Uber}. Nach \emph{Messungsfehler} und \emph{Repetitionsbeobachtungen} steht dort kein Doppelpunkt. Der zusätzliche Punkt in der Punktreihe bei
\emph{1822 Juli 10} und das kommaförmige Zeichen vor der
Überschrift \emph{B.} in der Verlagsanzeige sind wiedergegeben.
Ihre Funktion bleibt unklar. Die Verlagsanzeige folgt dem
Schlussblatt des Göttinger Exemplars und fehlt im Google-Digitalisat.
